% Methods: platform design, the five methods, the gate, and the autonomous
% per-commit generation methodology.

This section describes how the platform is built and how it was generated. It
covers four things in order: the five established methods that produce one daily
deliverable, the VVUQ gate and its thresholds, the Stage 1 automation that
surrounds the pipeline together with the Stage 2 deployment references, and the
autonomous per commit process that authored the codebase. Each claim is grounded
in the repository sources collected in \autoref{tab:methods_sources}; the prose
references those sources rather than reproducing them, and long module names are
abbreviated in the tables.

\begin{table}[H]
\caption{Methods source map (abbreviated paths under the repository root).}
\label{tab:methods_sources}
\centering
\begin{tabularx}{\textwidth}{L{4.8cm} >{\raggedright\arraybackslash}X}
\toprule
\textbf{Source (under root)} & \textbf{Use for} \\
\midrule
  \texttt{pipeline/} : 6 modules and \texttt{examples-pipeline/} & The four producing stages, the orchestrator, and the daily deliverable model. \\
  \texttt{vvuq/} : 4 modules and \texttt{examples-vvuq/} & The fifth method, the gate, and the three VVUQ dimensions. \\
  \texttt{configs/} : pipeline and vvuq YAML & The 200K cap, the 3 run default, and the gate thresholds. \\
  \texttt{physical-ai/} : factory and pilot modules & Stage 2 references and the safety interlocks. \\
  \texttt{simulation/}, \texttt{ingestion/}, \texttt{chunking/}, \texttt{scheduler/} & Stage 1 automation around the pipeline. \\
  \texttt{scripts/} verify script; \texttt{tests/} : 8 modules & The environment gate and the verification floor. \\
  \texttt{execution/01-foundation/} reports & Reproducibility: Python 3.11.15, tooling, clean static analysis. \\
\bottomrule
\end{tabularx}
\end{table}

\subsection{Five established methods, one daily deliverable}

The pipeline packages five methods that prior projects established and proved,
and runs them so that one daily deliverable passes through all five in a fixed
order. Four of the methods produce the deliverable and the fifth clears it. The
four producing stages are instruction generation, code generation, code
execution, and paper assembly, and they are realized in \texttt{pipeline/} by
\texttt{instruction\_stage.py}, \texttt{codegen\_stage.py},
\texttt{execution\_stage.py}, and \texttt{paper\_stage.py}. 

A single
data model in \texttt{deliverable.py} carries the work between stages: it defines
the stage enumeration, the per stage status, the artifact records, and the
\texttt{Deliverable} object itself, so that each stage reads the previous stage's
output and appends its own. The orchestrator in \texttt{orchestrator.py} threads
one \texttt{Deliverable} object through the four producing stages in strict order
and then invokes the fifth method, the VVUQ gate, as a callable applied to the
finished candidate. The orchestrator is configured with the number of simulation
runs, and the default configuration uses three.

Each producing stage emits a concrete artifact. Instruction generation writes a
bracketed instruction set, code generation writes a generated module, execution
writes a captured log whose key and value lines are parsed back into the
deliverable's metadata, and paper assembly writes a short paper that tabulates
the executed results. The acceleration that the pipeline reports is not a free
parameter. It follows directly from a single audited formula in the code
generation stage, \texttt{automated\_days = 30.0 / (1.0 + runs * 0.5)}, evaluated
against a fixed conventional baseline of 30 days. With the default of three
simulation runs the denominator is 2.5, the automated schedule is 12 days, and
the acceleration factor is the ratio of the two, namely 2.5. The same model
evaluated at one and two runs gives the intermediate factors of 1.5 and 2.0, so
the headline 2.5 is the third point on a transparent curve rather than an
isolated claim. Because the formula depends only on the run count and not on the
deliverable's topic, three identical inputs yield three identical outputs, a
property that matters when the uncertainty dimension is examined in
\autoref{sec:results}.

The four producing stages and their artifacts are summarized in
\autoref{tab:methods_stages}. Each stage is small, single purpose, and
deterministic in the absence of a live model backend, which is what allows the
whole pipeline to run in well under a millisecond and to be reproduced exactly.
The orchestrator does not hide this structure: it exposes each stage result and
each artifact on the deliverable, so the VVUQ gate that follows can inspect the
work rather than trust it. This separation, generation that produces and a gate
that inspects, is the architectural expression of the thesis, and it is why the
remainder of the methods concentrates on the gate.

\begin{table}[H]
\caption{The four producing stages, their modules, and their artifacts.}
\label{tab:methods_stages}
\centering
\begin{tabularx}{\textwidth}{L{2.7cm} L{3.6cm} L{3.0cm} >{\raggedright\arraybackslash}X}
\toprule
\textbf{Stage} & \textbf{Module} & \textbf{Artifact} & \textbf{Behavior} \\
\midrule
  Instruction & \texttt{instruction\_stage.py} & \texttt{instructions.md} & Emits the bracketed instruction set. \\
  Code generation & \texttt{codegen\_stage.py} & \texttt{generated\_*.py} & Authors code from the instructions. \\
  Code execution & \texttt{execution\_stage.py} & \texttt{execution\_log.txt} & Runs the code and parses key and value lines into metadata. \\
  Paper assembly & \texttt{paper\_stage.py} & \texttt{paper.md} & Assembles a short paper tabulating results. \\
\bottomrule
\end{tabularx}
\end{table}

\subsection{The VVUQ gate and its thresholds}

The fifth method is the part the thesis turns on. The gate in
\texttt{vvuq/vvuq\_gate.py} combines three independent checks, each implemented in
its own module, and decides whether a candidate deliverable may ship, as seen in \autoref{figwheel}. Its
thresholds are not hard coded into the gate; they are loaded at run time from
\texttt{configs/vvuq\_thresholds.yaml}, so the policy is auditable and can be
adjusted without modifying decision logic. Full set of defaults is in \autoref{tab:methods_thresholds}.

The verification module in \texttt{verification.py} answers whether the
deliverable was built correctly. It runs six structural and policy checks, namely
that all stages completed, that the expected artifacts are present, that those
artifacts are non empty, that no artifact exceeds the 200{,}000 byte per file
cap, that the schema is valid, and that the deliverable lints cleanly. The
fraction of these checks that pass must equal exactly 1.0; nothing less is
acceptable, so a single failed check among the six is sufficient to block. The
validation module in \texttt{validation.py} answers whether the right deliverable
was built. It compares the observed results against a trusted reference,
requiring an agreement of at least 0.95 and a maximum relative error of at most
0.05, and it additionally requires that a human review has been recorded. 

The
uncertainty module in \texttt{uncertainty.py} answers how confident the result
is. It computes the mean, the standard deviation, and the coefficient of
variation of each metric across the simulation runs, requires at least three
runs, and requires the coefficient of variation to stay at or below 0.10 at a
0.95 confidence level. The gate policy itself is deliberately strict: it blocks
on any failure of any dimension, and when the runs diverge beyond the dispersion
bound it does not silently fail but escalates the case to a human. This
escalation path is what separates a doubtful result, which a person should
adjudicate, from a clearly bad one, which the gate simply refuses. The posture is
consistent with the credibility and software life cycle expectations of medical
computational evidence~\cite{iec62304, iche6r3}.

\begin{table}[H]
\caption{VVUQ gate default thresholds, loaded verbatim from \texttt{configs/vvuq\_thresholds.yaml}.}
\label{tab:methods_thresholds}
\centering
\begin{tabularx}{\textwidth}{L{4.0cm} >{\raggedright\arraybackslash}X L{3.3cm}}
\toprule
\textbf{Dimension} & \textbf{Requirement} & \textbf{Default} \\
\midrule
  Verification & Fraction of structural checks passed & $= 1.0$ (6 of 6) \\
               & Schema valid                         & required \\
               & Lint clean                           & required \\
               & Per file size cap                    & 200{,}000 bytes \\
  \addlinespace
  Validation   & Agreement with reference             & $\geq 0.95$ \\
               & Maximum relative error               & $\leq 0.05$ \\
               & Recorded human review                & required \\
  \addlinespace
  Uncertainty  & Coefficient of variation across runs & $\leq 0.10$ \\
               & Consensus runs                       & $\geq 3$ \\
               & Confidence level                     & 0.95 \\
  \addlinespace
  Gate policy  & Block on any failure                 & enabled \\
               & Escalate on divergence               & enabled \\
\bottomrule
\end{tabularx}
\end{table}



Two design choices in the gate deserve emphasis because they are what make it
substantial rather than cosmetic. First, the verification fraction is an equality
to 1.0, not an inequality with slack. Many quality systems accept a high pass
rate, but a single unmet structural or policy check, for example an artifact one
byte over the cap, is treated here as disqualifying, so the gate cannot be
satisfied by being mostly correct. Second, the uncertainty module does not merely
report a coefficient of variation; it distinguishes a value that is high from one
that is divergent, and only divergence triggers the escalation path. The
practical effect is a three way outcome rather than a binary pass or fail: a clean
candidate is accepted, a clearly defective one is blocked, and an ambiguous one,
where the runs disagree enough that no single run can be trusted, is routed to a
human for adjudication. That third outcome is the part a purely automated gate
usually lacks, and it is the part that a regulated setting most needs.

\newpage

\vspace*{-1.25cm}

\begin{figure}[H]
    \hspace{-0.8cm}%
      {\includegraphics[width=1.1\linewidth, trim={0 25 0 0}, clip]{Images/04-vvuq-assurance-wheel.png}}
    \vspace{-0.5cm}
    \caption{VVUQ assurance wheel: the accepted deliverable clears every spoke, while blocked cases fail.}
    \label{figwheel}
\end{figure}

\subsection{Stage 1 automation and the Stage 2 deployment references}

Around the five methods sit four Stage 1 capabilities that make the pipeline run
continuously and safely. Triple simulation, in \texttt{simulation/}, runs each
project three times under distinct seeds through \texttt{triple\_runner.py} and
forms a consensus across the runs in \texttt{consensus.py}, flagging any metric
whose dispersion exceeds the bound as divergent. Robust ingestion, in
\texttt{ingestion/}, provides a web search client and a PDF processor that both
guard their optional dependencies: when a backend or the network is unavailable
the client retries up to the configured budget of four attempts and then degrades
to a clearly labeled offline stub rather than failing. 

Autochunking, in
\texttt{chunking/}, enforces the 200{,}000 byte per file cap by splitting
oversized files in \texttt{chunker.py} and writing a reconstruction README per
chunk in \texttt{readme\_generator.py}, so that any split file can be rebuilt with
a single concatenation. The non stop commit scheduler, in
\texttt{scheduler/commit\_scheduler.py}, plans an evenly spaced cadence of
deliverables across the day; the default configuration plans twenty four slots at
an hourly interval. 

These defaults, including the 200{,}000 byte cap, the three
run default, the four retry attempts, and the twenty four commits per day, live
in \texttt{configs/pipeline\_config.yaml}. Two further ingestion defaults sit in
the same configuration, a maximum of 64 documents per run and the four attempt
retry budget, and the chunker recognizes the project's text formats, the Python,
Markdown, LaTeX, JSON, and YAML sources, so that any of them can be split and
rebuilt. The reconstruction README that accompanies each split records the exact
concatenation command, so a chunked file is never left without the recipe that
restores it byte for byte.

The Stage 2 deployment references, in \texttt{physical-ai/}, show where the same
assurance posture is meant to extend. The lights off factory controller in
\texttt{lights\_off\_factory.py} runs autonomous batches across cells under a set
of safety interlocks, and the hybrid surgery and medicine pilot in
\texttt{hybrid\_surgery\_medicine.py} expresses the 2030 pancreatic cancer analog
as a timeline. Both are deliberately disabled by default in the deployment
configuration and require human oversight to run. Critically, one of the four
default factory interlocks is named \texttt{vvuq\_gate\_online}, alongside
\texttt{estop\_armed}, \texttt{perimeter\_clear}, and
\texttt{human\_oversight\_remote}. 

Wiring the Stage 1 assurance gate into the
Stage 2 safety chain as an interlock is the design decision that connects the two
stages: a physical batch cannot start unless the same gate that clears software
deliverables is itself online, which is essential when a 60 second procedure
leaves no time for manual verification inside the loop. The controller is a small
state machine. 

It idles in a dark state, transitions to running only when every
interlock is satisfied, returns to dark when a batch completes within its fault
budget, emergency stops when cumulative faults exceed that budget, and refuses to
start at all, raising a fault, when any interlock is unsatisfied. Because
\texttt{vvuq\_gate\_online} is one of those interlocks, a lapse in software
assurance is, by construction, a condition that prevents the physical line from
running. The clinical content of this pilot remains a planning reference and requires VVUQ clearance, human
oversight, IRB approval, and regulatory authorization before any real use.

\subsection{Autonomous per commit generation}

The platform was not written by hand. Claude Code Opus 4.7 (1M context) Max
generated the codebase autonomously in a managed cloud
container~\cite{anthropicopus47, anthropicclaudecodeweb}, and the way it did so
is itself part of the method, because it demonstrates that the assurance
discipline can be applied at the granularity of a single commit. The agent wrote
one section of the platform per commit and pushed each commit to GitHub in real
time, so that every increment was independently visible and the working context
stayed focused on one concern at a time. After the producing commits it re ran
the full continuous integration surface end to end and repaired any issue in a
dedicated second to last commit, then made a final commit that updated the
repository documentation. All commits belonged to a single pull request on
\texttt{cancer-automated}, and no other repository was touched~\cite{repo-cancer-automated}.
The seven commits that built the platform are summarized in
\autoref{tab:methods_codegen_commits}, where the gap column records the elapsed
time between consecutive commits and the churn column records the added and
removed lines.

\begin{table}[H]
\caption{Autonomous code generation commits that built the platform.}
\label{tab:methods_codegen_commits}
\centering
\begin{tabularx}{\textwidth}{C{0.7cm} L{2.0cm} >{\raggedright\arraybackslash}X C{1.6cm} C{2.0cm}}
\toprule
\textbf{\#} & \textbf{SHA} & \textbf{Section} & \textbf{Gap} & \textbf{Churn} \\
\midrule
  C1 & c9e1404 & Foundation: tooling, CI, community files & - & +1088 -2 \\
  C2 & a111ca8 & Pipeline (five established methods) & 4m42s & +825 \\
  C3 & 187ba9e & VVUQ framework & 2m44s & +691 \\
  C4 & dcb11c0 & Stage 1 automation (4 modules) & 4m46s & +1127 \\
  C5 & 082fa9f & Stage 2 deployment references & 1m49s & +389 \\
  C6 & 13b253d & Fix all errors (second to last) & 0m48s & +3 -1 \\
  C7 & bb27fd6 & Repository updates (last) & 3m05s & +119 -9 \\
\bottomrule
\end{tabularx}
\end{table}

The composition of the platform that these commits produced is summarized in two
compact inventories. \autoref{tab:methods_filetypes} breaks the 76 committed
files down by type, and \autoref{tab:methods_dirlines} gives the line count of
each source directory. Two observations from these inventories support the
thesis. First, the assurance and test code is comparable in size to the
producing code: the verification, validation, uncertainty, and gate modules
together, plus the eight test modules, account for a large share of the lines, so
the effort spent deciding whether deliverables are trustworthy is on par with the
effort spent generating them. 

Second, the per commit cadence in
\autoref{tab:methods_codegen_commits} is fine grained and fast, with most
sections separated by a few minutes, which contrasts with prior non incremental
systems that generate an entire codebase in one batch and verify only at the end.
A representative prior run by the author generated 108 modules in a single batch
of thirteen commits~\cite{kawchak2026sponsor}; the present run instead gates each
section as it lands, so a regression is caught at the commit that introduced it
rather than at a terminal verification pass.

\begin{table}[H]
\caption{Committed files by type for the generated platform (76 files total).}
\label{tab:methods_filetypes}
\centering
\begin{tabularx}{\textwidth}{>{\raggedright\arraybackslash}X C{3.0cm} C{3.0cm}}
\toprule
\textbf{File type} & \textbf{Files} & \textbf{Lines} \\
\midrule
  Python (\texttt{.py})            & 44 & 2829 \\
  Markdown (\texttt{.md})          & 24 & 1140 \\
  YAML (\texttt{.yaml})            & 2  & 79 \\
  Workflow (\texttt{.yml})         & 1  & 73 \\
  Text (\texttt{.txt})             & 1  & 43 \\
  Config (\texttt{.toml})          & 1  & 32 \\
  Citation (\texttt{.cff})         & 1  & 30 \\
  LICENSE                          & 1  & 21 \\
  \texttt{.gitignore}              & 1  & 6 \\
  \midrule
  Total                            & 76 & 4253 \\
\bottomrule
\end{tabularx}
\end{table}

\begin{table}[H]
\caption{Source directory line counts for the generated platform.}
\label{tab:methods_dirlines}
\centering
\begin{tabularx}{\textwidth}{L{3.2cm} C{2.4cm} >{\raggedright\arraybackslash}X}
\toprule
\textbf{Directory} & \textbf{Lines} & \textbf{Composition} \\
\midrule
  \texttt{pipeline/}    & 753 & six stage modules plus three examples \\
  \texttt{vvuq/}        & 573 & four assurance modules plus three examples \\
  \texttt{tests/}       & 559 & eight test modules, 51 tests, plus conftest \\
  \texttt{physical-ai/} & 325 & factory and pilot modules plus two examples \\
  \texttt{ingestion/}   & 278 & web search and PDF processor plus examples \\
  \texttt{simulation/}  & 266 & triple runner and consensus plus examples \\
  \texttt{chunking/}    & 260 & chunker and README generator plus examples \\
  \texttt{.github/}     & 185 & CI workflow and community templates \\
  \texttt{scheduler/}   & 143 & commit scheduler plus one example \\
  \texttt{scripts/}     & 112 & environment verification script \\
  \texttt{configs/}     & 79  & pipeline and VVUQ threshold configuration \\
\bottomrule
\end{tabularx}
\end{table}

Reproducibility was treated as part of the method rather than an afterthought. The
environment was pinned to Python 3.11.15 and the suite was run under pytest
9.0.3, with the static analysis performed by the same \texttt{ruff} and
\texttt{yamllint} tools the continuous integration uses. A standalone verification
script confirms the environment before any deliverable runs, and the core engine
depends only on the standard library, with every heavy or optional dependency
imported behind a guard so that the package stays importable when a dependency is
absent. 

This is why the lint and test surface stays green on a minimal runner and
why the recorded commands can be replayed to reproduce each claim. The discipline
extends the prior humanoid instruction to code to execution to paper workflow by
adding exactly this assurance scaffolding around it~\cite{anthropicclaudecodeweb}. Taken together, the architecture and the build process embody the thesis in their
proportions. The generation is fast and inexpensive, the gate is strict and
auditable, and the assurance work, not the generation, dominated both the code
that was written and the discipline with which it was committed.
